% Options for packages loaded elsewhere
% Options for packages loaded elsewhere
\PassOptionsToPackage{unicode}{hyperref}
\PassOptionsToPackage{hyphens}{url}
\PassOptionsToPackage{dvipsnames,svgnames,x11names}{xcolor}
%
\documentclass[
  11pt,
]{article}
\usepackage{xcolor}
\usepackage[top=0.4in,bottom=0.45in,left=0.65in,right=0.65in,includehead,headheight=52pt,headsep=7pt,includefoot,footskip=24pt]{geometry}
\usepackage{amsmath,amssymb}
\setcounter{secnumdepth}{-\maxdimen} % remove section numbering
\usepackage{iftex}
\ifPDFTeX
  \usepackage[T1]{fontenc}
  \usepackage[utf8]{inputenc}
  \usepackage{textcomp} % provide euro and other symbols
\else % if luatex or xetex
  \usepackage{unicode-math} % this also loads fontspec
  \defaultfontfeatures{Scale=MatchLowercase}
  \defaultfontfeatures[\rmfamily]{Ligatures=TeX,Scale=1}
\fi
\usepackage{lmodern}
\ifPDFTeX\else
  % xetex/luatex font selection
\fi
% Use upquote if available, for straight quotes in verbatim environments
\IfFileExists{upquote.sty}{\usepackage{upquote}}{}
\IfFileExists{microtype.sty}{% use microtype if available
  \usepackage[]{microtype}
  \UseMicrotypeSet[protrusion]{basicmath} % disable protrusion for tt fonts
}{}
\makeatletter
\@ifundefined{KOMAClassName}{% if non-KOMA class
  \IfFileExists{parskip.sty}{%
    \usepackage{parskip}
  }{% else
    \setlength{\parindent}{0pt}
    \setlength{\parskip}{6pt plus 2pt minus 1pt}}
}{% if KOMA class
  \KOMAoptions{parskip=half}}
\makeatother
% Make \paragraph and \subparagraph free-standing
\makeatletter
\ifx\paragraph\undefined\else
  \let\oldparagraph\paragraph
  \renewcommand{\paragraph}{
    \@ifstar
      \xxxParagraphStar
      \xxxParagraphNoStar
  }
  \newcommand{\xxxParagraphStar}[1]{\oldparagraph*{#1}\mbox{}}
  \newcommand{\xxxParagraphNoStar}[1]{\oldparagraph{#1}\mbox{}}
\fi
\ifx\subparagraph\undefined\else
  \let\oldsubparagraph\subparagraph
  \renewcommand{\subparagraph}{
    \@ifstar
      \xxxSubParagraphStar
      \xxxSubParagraphNoStar
  }
  \newcommand{\xxxSubParagraphStar}[1]{\oldsubparagraph*{#1}\mbox{}}
  \newcommand{\xxxSubParagraphNoStar}[1]{\oldsubparagraph{#1}\mbox{}}
\fi
\makeatother


\usepackage{longtable,booktabs,array}
\newcounter{none} % for unnumbered tables
\usepackage{calc} % for calculating minipage widths
% Correct order of tables after \paragraph or \subparagraph
\usepackage{etoolbox}
\makeatletter
\patchcmd\longtable{\par}{\if@noskipsec\mbox{}\fi\par}{}{}
\makeatother
% Allow footnotes in longtable head/foot
\IfFileExists{footnotehyper.sty}{\usepackage{footnotehyper}}{\usepackage{footnote}}
\makesavenoteenv{longtable}
\usepackage{graphicx}
\makeatletter
\newsavebox\pandoc@box
\newcommand*\pandocbounded[1]{% scales image to fit in text height/width
  \sbox\pandoc@box{#1}%
  \Gscale@div\@tempa{\textheight}{\dimexpr\ht\pandoc@box+\dp\pandoc@box\relax}%
  \Gscale@div\@tempb{\linewidth}{\wd\pandoc@box}%
  \ifdim\@tempb\p@<\@tempa\p@\let\@tempa\@tempb\fi% select the smaller of both
  \ifdim\@tempa\p@<\p@\scalebox{\@tempa}{\usebox\pandoc@box}%
  \else\usebox{\pandoc@box}%
  \fi%
}
% Set default figure placement to htbp
\def\fps@figure{htbp}
\makeatother





\setlength{\emergencystretch}{3em} % prevent overfull lines

\providecommand{\tightlist}{%
  \setlength{\itemsep}{0pt}\setlength{\parskip}{0pt}}



 


% =====================================================================
%  _style.tex  ·  Design system for EIL research highlights (v4 layout)
%  Edit here to restyle ALL research highlights. Included via the .qmd YAML.
%  Matches the canonical single-column v4 template; the masthead and footer
%  are set per-document via fancyhdr in each .qmd's include-in-header.
% =====================================================================

% ---- Colours (exact hex) --------------------------------------------
\definecolor{ink}{HTML}{1D2620}        % headings, title, section heads, rule
\definecolor{body}{HTML}{33352C}       % main body paragraphs
\definecolor{accentred}{HTML}{641111}  % RESEARCH HIGHLIGHT label, links, fig title
\definecolor{boxred}{HTML}{8C2826}     % Bottom Line box fill
\definecolor{boxtext}{HTML}{EEF1EA}    % Bottom Line body text (off-white)
\definecolor{muted}{HTML}{50534A}      % figure source line
\definecolor{cite}{HTML}{6E6D61}       % citation line
\definecolor{faint}{HTML}{9A978A}      % footer
\definecolor{warmrule}{HTML}{D8D2C2}   % thin warm-grey dividers
\definecolor{paper}{HTML}{FAF8F1}      % warm page background

% \pagecolor{paper}                    % uncomment for warm off-white background

% ---- Fonts (Source Sans 3 / Source Serif 4) -------------------------
\usepackage[T1]{fontenc}
\usepackage[default]{sourcesanspro}
\usepackage{sourceserifpro}
\renewcommand{\familydefault}{\sfdefault}

% ---- Page -----------------------------------------------------------
\pagestyle{empty}
\setlength{\parindent}{0pt}
\setlength{\parskip}{2pt}
% Default body size/colour. The v4 .qmd overrides this to 10.5/13.5pt.
\AtBeginDocument{\color{body}\fontsize{9.4pt}{13.6pt}\selectfont}

% ---- Packages -------------------------------------------------------
\usepackage{microtype}
\usepackage{soul}
\setul{2pt}{.5pt}
\usepackage{tcolorbox}
\tcbuselibrary{skins}
\usepackage[colorlinks=true, urlcolor=accentred, linkcolor=accentred, pdfborder={0 0 0}]{hyperref}

% =====================================================================
%  MACROS
% =====================================================================

% --- Title + citation (citation has a thin warm divider beneath) ------
\newcommand{\brieftitle}[1]{%
  {\color{ink}\fontsize{18pt}{19.5pt}\selectfont\bfseries\textls[-12]{#1}\par}\vspace{3pt}}
\newcommand{\briefcitation}[1]{%
  {\color{cite}\fontsize{9.4pt}{12pt}\selectfont #1\par}\vspace{3pt}%
  \noindent\textcolor{warmrule}{\rule{\linewidth}{0.7pt}}}

% --- Section heads (uppercase, ink) ----------------------------------
\newcommand{\sectionhead}[1]{%
  \vspace{2pt}{\color{ink}\bfseries\fontsize{9.75pt}{12pt}\selectfont
  \textls[20]{\MakeUppercase{#1}}\par}\vspace{1pt}}

% --- Sub-heads (small uppercase) — Takeaway subheads -----------------
% The v4 .qmd overrides this to the brand maroon via \renewcommand.
\newcommand{\subhead}[1]{%
  {\color{ink}\bfseries\fontsize{8.25pt}{10pt}\selectfont
  \textls[50]{\MakeUppercase{#1}}\par}\vspace{2pt}}

% --- Figure title + source line --------------------------------------
\newcommand{\figtitle}[2]{%
  {\bfseries\color{accentred}\fontsize{7.1pt}{9pt}\selectfont\textls[40]{\MakeUppercase{#1}}}\par%
  {\color{muted}\fontsize{7.1pt}{9pt}\selectfont #2}\par\vspace{4pt}}

% =====================================================================
%  COLOURED BOXES
% =====================================================================

% --- Bottom Line box (border-radius 6px -> arc=4pt) ------------------
\newtcolorbox{bottomline}{%
  enhanced, arc=4pt, outer arc=4pt, boxrule=0pt,
  colback=boxred, coltext=boxtext,
  left=7pt, right=7pt, top=7pt, bottom=7pt,
  before upper={\fontsize{9.4pt}{13.6pt}\selectfont}}

% --- Info box (contact + full citation; centered, thin faint border) --
% Superseded by the moreinfo environment below; kept so older highlights
% still compile. Faint text at footer size. Wrap any link as
% \href{url}{\color{faint}\underline{text}} so it inherits the faint
% colour instead of the global accent link colour.
\newtcolorbox{infobox}{%
  enhanced, width=0.9\linewidth, center, arc=4pt, outer arc=4pt,
  boxrule=0.5pt, colframe=faint, colback=white, coltext=faint,
  left=9pt, right=9pt, top=7pt, bottom=7pt,
  before upper={\fontsize{7.4pt}{9.5pt}\selectfont}}

% =====================================================================
%  CLOSING BLOCK
% =====================================================================

% --- For more information (full width; rule above, no box) ------------
% Citation, author affiliations + contact, then the EIL boilerplate —
% one paragraph each. Wrap any link as
% \href{url}{\color{cite}\underline{text}} so it inherits the block's
% grey instead of the global accent link colour.
\newenvironment{moreinfo}{%
  \par\vspace{12pt}%
  \noindent\textcolor{warmrule}{\rule{\linewidth}{0.7pt}}\par
  \vspace{7pt}%
  {\color{ink}\bfseries\fontsize{8.25pt}{10pt}\selectfont
   \textls[50]{\MakeUppercase{For more information}}\par}\vspace{4pt}%
  \color{cite}\fontsize{8.4pt}{11.5pt}\selectfont
  \setlength{\parskip}{5pt}%
}{\par}
% --- Running masthead + footer (logo, label, rule, page number) ---
\usepackage{fancyhdr}
\renewcommand{\headrulewidth}{0pt}
\renewcommand{\footrulewidth}{0pt}
\fancypagestyle{brief}{%
  \fancyhf{}%
  \fancyhead[L]{%
    \begin{minipage}[c]{0.62\textwidth}\raggedright\color{ink}
\includegraphics[height=0.45in]{../../../formats/logos/eil-logo-maroon.png}\end{minipage}\hfill%
    \begin{minipage}[c]{0.36\textwidth}\raggedleft\fontsize{8.25pt}{10pt}\selectfont\bfseries\color{accentred}\textls[200]{\MakeUppercase{Research Highlight}}\end{minipage}%
    \par\vspace{2pt}%
    \noindent\textcolor{ink}{\rule{\textwidth}{2.2pt}}%
  }%
  \fancyfoot[C]{%
    \color{warmrule}\rule{\textwidth}{0.7pt}\par\vspace{3pt}%
    \hbox to\textwidth{%
      \color{faint}\fontsize{7.4pt}{9pt}\selectfont
      For full paper:\hspace{0.4em}\href{https://researchonline.lse.ac.uk/id/eprint/126795}{\color{faint}\underline{\textls[40]{researchonline.lse.ac.uk/id/eprint/126795}}}%
      \hfill
      Page \thepage%
    }%
  }%
}
\pagestyle{brief}
% Subhead macro retained from the design system (unused here, since the
% Takeaways subheads give way to a flowing concluding section).
\renewcommand{\subhead}[1]{%
  {\color{accentred}\bfseries\fontsize{8.25pt}{10pt}\selectfont
  \textls[50]{\MakeUppercase{#1}}\par}\vspace{2pt}}
\makeatletter
\@ifpackageloaded{caption}{}{\usepackage{caption}}
\AtBeginDocument{%
\ifdefined\contentsname
  \renewcommand*\contentsname{Table of contents}
\else
  \newcommand\contentsname{Table of contents}
\fi
\ifdefined\listfigurename
  \renewcommand*\listfigurename{List of Figures}
\else
  \newcommand\listfigurename{List of Figures}
\fi
\ifdefined\listtablename
  \renewcommand*\listtablename{List of Tables}
\else
  \newcommand\listtablename{List of Tables}
\fi
\ifdefined\figurename
  \renewcommand*\figurename{Figure}
\else
  \newcommand\figurename{Figure}
\fi
\ifdefined\tablename
  \renewcommand*\tablename{Table}
\else
  \newcommand\tablename{Table}
\fi
}
\@ifpackageloaded{float}{}{\usepackage{float}}
\floatstyle{ruled}
\@ifundefined{c@chapter}{\newfloat{codelisting}{h}{lop}}{\newfloat{codelisting}{h}{lop}[chapter]}
\floatname{codelisting}{Listing}
\newcommand*\listoflistings{\listof{codelisting}{List of Listings}}
\makeatother
\makeatletter
\makeatother
\makeatletter
\@ifpackageloaded{caption}{}{\usepackage{caption}}
\@ifpackageloaded{subcaption}{}{\usepackage{subcaption}}
\makeatother
\usepackage{bookmark}
\IfFileExists{xurl.sty}{\usepackage{xurl}}{} % add URL line breaks if available
\urlstyle{same}
\hypersetup{
  colorlinks=true,
  linkcolor={blue},
  filecolor={Maroon},
  citecolor={Blue},
  urlcolor={Blue},
  pdfcreator={LaTeX via pandoc}}


\author{}
\date{}
\begin{document}


% ---------------------------------------------------------------------
% TITLE + CITATION  (masthead repeats in the page header)
% ---------------------------------------------------------------------
\brieftitle{Coal's Decline Imposed Large, Lasting Costs on Workers}
\briefcitation{Lessons from ``Transitional Costs and the Decline of Coal: Worker-Level Evidence'' by Jonathan Colmer, Eleanor Krause, Eva Lyubich \& John Voorheis, Working Paper, 2026.}

\vspace{-2pt}

% override body font size for this document
\fontsize{10.5pt}{13.5pt}\selectfont

% ---------------------------------------------------------------------
% BOTTOM LINE  (~40–50 words)
% ---------------------------------------------------------------------
\begin{bottomline}
{\fontsize{11pt}{15pt}\selectfont
\textbf{When U.S. coal production and employment suddenly fell by half in the 2010s, the workers it left behind never caught back up.} The most coal-dependent workers lost close to two years' worth of earnings over the following decade --- and neither switching industries nor moving to a new region did much to help recoup those losses.
}
\end{bottomline}

\vspace{6pt}

% ---------------------------------------------------------------------
% BACKGROUND  (broad — poses the question)
% ---------------------------------------------------------------------
\sectionhead{Background}
\noindent U.S. coal mining employment has fallen more than half since its peak in 2011, as cheap natural gas from the fracking boom and falling renewable energy costs pushed coal out of the power sector. Earlier research shows that coal \textit{communities} lost jobs, earnings, and tax revenue. However, aggregate numbers may hide how individual workers fared. Some may have quickly moved to new jobs, while others may have borne heavy \textbf{transitional costs}: the lasting losses to earnings and employment that workers bear when an industry shrinks. Colmer et al. wanted to use the decline of coal to understand individual-level transitional costs, and how that might inform the broader question of what the energy transition could mean for workers in other high-wage, regionally concentrated industries.

\noindent 
To assess the private costs of coal's decline and understand if aggregated patterns held for individual workers, Colmer et al. assembled administrative records covering nearly every person who worked in U.S. coal mining before, during, and after coal's 2011 peak. This let them follow each worker's earnings, employment, and location year by year, even if workers worked outside coal or were unemployed. They could then compare those paths with otherwise similar workers who depended less on coal for their livelihoods. \textbf{This comparison showed that coal-dependent workers experienced persistent losses in earnings and employment following coal's decline.}

\vspace{6pt}

% ---------------------------------------------------------------------
% FIGURE — recreation of the paper's Fig. 7(a); built by code/make-fig1.R
%   Placed after the background so it lands on page 1: the bottom-line box
%   has already stated the loss, so the chart illustrates a claim the
%   reader holds rather than spoiling one they don't.
% ---------------------------------------------------------------------
\begin{center}
\begin{minipage}{0.82\linewidth}
\figtitle{Coal workers' earnings fell further behind comparable workers every year}{Source: recreated from Colmer et al.\ (2026), Fig.\ 7a.}
% Scaled to 0.94 so the figure block clears the page-1 text block; at full
% width it overruns by ~2pt and the unbreakable minipage jumps to page 2.
\centerline{\includegraphics[width=0.94\linewidth]{../data-viz/figures/fig1.png}}
\par\vspace{3pt}
{\color{faint}\fontsize{7.1pt}{9.5pt}\selectfont \textit{Note:} Each point is the cumulative earnings gap between workers whose livelihoods depended most on coal and otherwise similar workers who depended on it less, measured in years of pre-decline pay. After coal's 2011 peak, more- and less-coal-dependent workers' earnings steadily diverge --- by 2019 the gap reaches 1.8 years' worth of earnings. Values digitized from the published figure.\par}
\end{minipage}
\end{center}

\vspace{6pt}

% ---------------------------------------------------------------------
% HOW THE STUDY TESTED IT  (narrowing — data + research design)
% ---------------------------------------------------------------------
\sectionhead{How the study tested it}
\noindent Measuring what coal's decline cost workers is challenging, because a career is shaped by many forces at once. One might simply track coal workers' earnings after 2011 and call the drop the cost --- but aging, the broader economy, and ordinary job churn were all pulling on those same paychecks, so a raw decline can't be pinned on coal alone. The next step is to compare coal workers against everyone else. Yet coal workers tend to be better paid, clustered in a few isolated regions, and different from other workers in age and education, so any earnings gap between them could reflect those existing differences rather than what coal's decline did to them.

\noindent Colmer et al. narrowed the comparison to workers who looked alike on those traits --- evening out age, education, earnings, and location --- and set the most coal-dependent workers beside otherwise similar workers who depended on coal far less. This holds only if the two groups would have followed the same earnings path absent the industry's collapse, which the researchers have to assume. That assumption is easier to credit here because the collapse was driven by national energy prices rather than any one worker's choices, so the most coal-dependent workers were hit by a force outside their control. If coal's decline had left workers unharmed, or affected both groups equally, no earnings gap should have opened between them.

\vspace{6pt}

% ---------------------------------------------------------------------
% WHAT THE STUDY FOUND  (the result, then why workers couldn't recover)
% ---------------------------------------------------------------------
\sectionhead{What the study found}
\noindent This is not what the researchers found. Between 2012 and 2019, \textbf{the most coal-dependent workers lost cumulative earnings equal to roughly 1.8 times a full pre-decline year of pay}. The losses came through two channels at once: compared with similar workers, coal-dependent workers spent about five more months out of work \textit{and}, in the years they did work, earned roughly 19 percent less.

\noindent Workers couldn't simply move or retrain their way out. Coal paid well --- about \$13,000 a year more than comparable work --- and its jobs are concentrated in a few isolated regions where entire economies lean on the industry, so nearby alternatives were scarce. \textbf{Those who left coal for other industries earned over 30 percent less than they had in coal, and relocating to a new labor market barely helped.} Where workers displaced by factory closings often find their way into new jobs, a large share of coal workers stopped working altogether, a sign their skills didn't transfer to the jobs around them.

\noindent As earnings fell, coal-dependent workers leaned on government transfers, especially Social Security Disability Insurance (SSDI), and drew from retirement savings earlier. These payments cushioned the blow but replaced only a fraction of what workers lost. \textbf{A decade on, coal's decline had carved a deep and persistent hole in workers' earnings that neither the labor market nor the safety net filled.}

\noindent Coal's decline is an early look at what shrinking demand for fossil fuels could mean for workers in regionally concentrated, high-wage industries. Coal is just one industry, and its geography is unusually isolating, so its exact path won't repeat everywhere. What the study establishes is narrower and still sobering: \textbf{when a high-wage, regionally concentrated industry contracts, its workers can bear the cost for a decade, and existing safety nets catch only part of the fall.}

\vspace{8pt}

\begin{moreinfo}
Colmer, Jonathan, Eleanor Krause, Eva Lyubich, and John Voorheis. 2026. ``Transitional Costs and the Decline of Coal: Worker-Level Evidence.'' Working paper. \href{https://researchonline.lse.ac.uk/id/eprint/126795}{\color{cite}\underline{researchonline.lse.ac.uk/id/eprint/126795}}.\par

Jonathan Colmer is at the University of Virginia. Eleanor Krause is at the University of Kentucky. Eva Lyubich and John Voorheis are at the US Census Bureau. Contact Eleanor Krause (eleanor.krause@uky.edu) with questions about this research.\par

The \href{https://www.environmental-inequality-lab.org/}{\color{cite}\underline{Environmental Inequality Lab}} is a nonpartisan research group that applies a rigorous, data-driven approach to understand how our environment shapes economic opportunity and well-being.\par
\end{moreinfo}




\end{document}
